% Cote 119: offprints, reports, courses, CV and letters (typed and printed part): archive cote 119, site page 238 (PDF page 239).
% typescript. Draft compared: none.
% AI-transcribed by Claude Opus 5.5 (claude-opus-5-5, adaptive thinking, effort high), Message Batches API, batch r119_full (msgbatch_01DwZNm5Y7YsKBWyAN4sFs8w), 2026-11-13; not yet checked by eye.
\scanpage{119/238}
\mspage{- 26 -}

M. DEMAZURE, Motifs des variétés algébriques,
Sém. Bourbaki n° 365, 1969/70.

A. DOUADY, Le Problème des modules pour les sous-%
espaces analytiques compacts..., Ann.
Inst. Fourier, vol. 16, pp. 1-98 (1966).

O. FORSTER et K. KNORR, Relativ-analytische Räume
und die Kohärenz von Bildgarben, Inventiones
Math. Vol. 16, pp. 113-160 (1972).

I.M. GELFAND et N. Ja. VILENKIN, Les distributions,
tome 4, Applications de l'analyse harmonique,
Dunod, Paris, 1968 (traduction).

J. GIRAUD, Cohomologie non abélienne, Grundlehren
des Math. Wiss. Bd. 179, 1971, Springer.

M. HAKIM, Topos annelés et schémas relatifs, Ergebnisse
des Math. Bd. 64, 1972, Springer.

R. HARTSHORNE, Residues and Duality, Lecture Notes
in Mathematics n° 20 (1966).

L. ILLUSIE, Complexe Cotangent et Déformations I,
Lecture Notes in Math. n° 239 (1971),
Springer et II, idem, n° 283 (1972).

R. KIEHL, Relativ analytische Raume, Inventiones
Math. vol. 16, pp. 40-112 (1972).

D. KNUTSON, Algebraic Spaces, Lecture Notes in Math.
n° 203 (1971), Springer.

F.N. LAWVERE, Toposes, algebraic geometry and logic,
Lecture notes in Math., n° 274 (1972), Springer.

I. MANIN, a) Théorie des groupes formels commutatifs
sur des corps de caractéristique finie,
(en russe) Uspekhi mat. Nauk, 1963, t. 18,
pp. 3-90. (Il existe une traduction anglaise
de l'Amer. Math. Soc.).
b) Correspondances, motifs et transformations
monoïdales (en russe), Mat. Sbornik t. 77,
pp. 475-507.

\margin{\unsure{1,238}}
